\section{Computational certificates and formal verification}\label{sec:verification}
The positive coloring certificates are checked using integer arithmetic and
the coordinate definition of orthogonality. The full checker imports neither
the search program nor a SAT solver. It establishes enumeration completeness
from distinctness and Gaussian counts, then checks every pair of vertices.
Three deliberately invalid certificates, respectively omitting a vertex,
duplicating a subspace, and identifying the colors of adjacent coordinate
lines, are rejected. The line-coloring checker independently verifies its
partition and all within-class pairs.

The radical/quotient clique theorem has an existing Lean formalization.
Its source is \path{formal/NikandishClique.lean}, with theorem
\texttt{NikandishClique.result} in its documented namespace. The exact
upstream revision, Lean version and dependencies are recorded in
\texttt{formal/UPSTREAM.json}. The new chromatic certificate is currently
verified by the independent integer checker; its Lean verification is being
developed separately from the archived clique theorem.

The public package~\cite{Package} contains the literal witnesses, independent
checkers, search records, and SHA-256 bindings. The new geometric profiler
checks every assigned color against its $H(U)$-determined list. The
dimension-seven reduction is also replayed in dimension six, where it
reproduces the established core size and Boolean encoding counts.
